\begin{lemma}
Let $g\in\mathsf{IncSeq}$ be nonidentity.  There is a prefix-continuous map
$\sigma:\mathsf{IncSeq}\to\mathsf{IncSeq}$ satisfying
\[
\sigma(f\circ\mathsf{SuccSeq})=\sigma(f)\circ g
\]
for every $f$.  Equivalently, the ordinary shift continuously
homomorphically maps to right composition by $g$.
\end{lemma}
\begin{proof}
By $hg$ and \noderef{FirstMovedPoint}, choose $k_g$ such that
$k_g<g(k_g)$ and $g(j)=j$ for every $j<k_g$.  Put
$G(n)=g^{[n]}(k_g)$.  For each $f$, choose an increasing injection
$\sigma(f)$ whose values are
$\sigma(f)(l)=\mathsf{BlockSigma}(g,k_g,f,l)$; this choice is supplied by
\noderef{BlockSigmaEmbedding}.  Since the chosen values are exactly the
block formula of \noderef{BlockSigma}, \noderef{BlockSigmaContinuous}
implies the following prefix condition: for every $f$ and output length
$n$, there is an input length $m$ such that every $h$ agreeing with $f$ on
the first $m$ coordinates satisfies
$\sigma(h)(l)=\sigma(f)(l)$ for all $l<n$.  Thus $f\mapsto\sigma(f)$ is
prefix-continuous, and its values lie in $\mathsf{IncSeq}$.

By the definitions of continuous morphism and its actions in
\noderef{ContMor}, \noderef{ShiftMap}, and \noderef{RightComp}, it remains
to verify the commuting equation coordinatewise.  Applying
\noderef{BlockSigmaIntertwines} with the witness $k_g$ gives
\[
\sigma_{f\circ\mathsf{SuccSeq}}(l)=\sigma_f(g(l)).
\]
The left side is $\sigma(\mathsf{ShiftMap}(f))(l)$ by
\noderef{SuccSeq}, and the right side is
$(\mathsf{RightComp}\ g\ (\sigma(f)))(l)$ by \noderef{RightComp}.
Therefore $\sigma(\mathsf{ShiftMap}(f))=\mathsf{RightComp}\ g\ (\sigma(f))$
for every $f$, so the prefix-continuous map $\sigma$ is the required
continuous homomorphism and establishes $\mathsf{ContMor}$.
\end{proof}
